\section{Oggetti e Array: Strutture Dati (Lezione 05)}

\subsection{Gli Oggetti (Objects)}
In JavaScript, gli oggetti sono contenitori per valori con etichetta (proprietà).
Pensa a un oggetto come a una "cosa" della vita reale. Una macchina ha proprietà come: peso, colore, modello.

\begin{concept}{Struttura: Chiave $\to$ Valore}
Pensa a un oggetto come a una "cosa" della vita reale. Una macchina ha proprietà come: peso, colore, modello.
\end{concept}
\begin{jsbox}
const auto = {
    marca: "Fiat",
    modello: "500",
    colore: "bianco"
};
\end{jsbox}

\subsubsection{Come accedere ai dati}
Ci sono due modi per leggere una proprietà:
\begin{enumerate}
    \item \textbf{Dot Notation (Punto):} \texttt{auto.marca} // "Fiat" (Più comune)
    \item \textbf{Bracket Notation (Parentesi):} \texttt{auto["marca"]} // "Fiat" (Utile se il nome della proprietà è in una variabile)
\end{enumerate}

\subsection{Memory: Valori vs Riferimenti}
Capire come varia la memoria è fondamentale. In JavaScript ci sono due regole diverse.

\subsubsection{1. Tipi Primitivi (Valori)}
\begin{concept}{Copia per Valore}
Numeri, Stringhe, Booleani. La variabile contiene \textbf{direttamente} il dato. Se copi, duplichi il dato.
\end{concept}
\begin{jsbox}
let a = 5;
let b = 5;
let c = a; // Copio il valore 5

console.log(a == b); // TRUE (5 e' uguale a 5)
console.log(a == c); // TRUE
\end{jsbox}

\subsubsection{2. Tipi Complessi (Riferimenti)}
\begin{concept}{Copia per Riferimento (Il Telecomando)}
Oggetti, Array, Funzioni.
La variabile contiene un \textbf{indirizzo di memoria} (un puntatore), non l'oggetto vero e proprio.
Se copi la variabile, copi solo il puntatore (telecomando), non l'oggetto (TV).
\end{concept}
\begin{jsbox}
let list1 = [1, 2, 3];
let list2 = [1, 2, 3];

// Sembrano uguali, ma sono due casi diversi in memoria!
console.log(list1 == list2); // FALSE (Indirizzi diversi)

let list3 = list1; // Copio l'indirizzo (il telecomando)
console.log(list1 == list3); // TRUE (Puntano allo stesso array)
\end{jsbox}

\subsubsection{Modifica Condivisa}
\begin{concept}{Side Effect (Effetto Collaterale)}
Poiché \texttt{list3} e \texttt{list1} puntano allo stesso oggetto, se ne modifico uno, cambia anche l'altro.
\end{concept}
\begin{jsbox}
list3.push(99);
console.log(list1); // [1, 2, 3, 99] -> VEDO LA MODIFICA!
\end{jsbox}

\begin{concept}{Visualizzazione Memoria: Stack vs Heap}
\begin{itemize}
    \item \textbf{Stack:} Veloce, per numeri e boolean (Valori).
    \item \textbf{Heap:} Grande e dinamico, per Oggetti e Array (Riferimenti).
\end{itemize}
\end{concept}

\begin{center}
\begin{tikzpicture}[
    node distance=0cm,
    stacknode/.style={draw, rectangle, minimum width=2.5cm, minimum height=0.8cm, fill=white},
    heapnode/.style={draw, rectangle, rounded corners, minimum width=3cm, minimum height=1.5cm, fill=niceBlue!10, align=center},
    label/.style={font=\bfseries\small, color=gray}
]
% ... (TikZ Code remained same, simplified here for brevity but assuming user wants it kept) ...
% Stack
\node[label] at (0, 3.5) {STACK (Esecuzione)};
\node[stacknode] (var1) at (0, 2.5) {age = 25};
\node[stacknode] (var2) at (0, 1.7) {name = "Luca"};
\node[stacknode, fill=jsYellow!20] (ref) at (0, 0.9) {user = <Ref \#01>};

% Heap
\node[label] at (6, 3.5) {HEAP (Memoria Oggetti)};
\node[heapnode] (obj) at (6, 1) {
    \textbf{Object \#01} \\
    \{ \\
    name: "Marco", \\
    active: true \\
    \}
};

% Puntatore
\draw[-{Stealth[length=3mm]}, thick, color=niceBlue] (ref.east) to[out=0, in=180] node[midway, above, font=\tiny\ttfamily] {Puntatore} (obj.west);
\end{tikzpicture}
\end{center}

\subsection{Visualizzazione Array}
\begin{center}
\begin{tikzpicture}
    % Definizione stile matrice array
    \matrix (A) [matrix of nodes, nodes={draw, minimum size=1cm, anchor=center, fill=white}, column sep=-\pgflinewidth]
    { 10 & 45 & 22 & 99 & 7 \\ };
    
    % Indici
    \foreach \i [count=\xi from 0] in {1,...,5}
        \node[above=0.1cm of A-1-\i, font=\ttfamily\tiny, color=gray] {\xi};
        
    % Etichetta
    \node[left=0.5cm of A, font=\bfseries] {const arr = };
    
    % Freccia puntatore
    \draw[<-] (A-1-3.south) -- ++(0, -0.5) node[below, font=\small] {arr[2]};
\end{tikzpicture}
\end{center}

\subsection{Array (Vettori)}
\begin{concept}{Lista Ordinata}
Un array è una variabile speciale che può contenere più valori in una lista ordinata $(0, 1, 2, \dots)$.
\end{concept}
\begin{jsbox}
const frutta = ["Mela", "Banana", "Pera"];
console.log(frutta[0]); // "Mela"
\end{jsbox}

\subsubsection{Aggiungere e Rimuovere Elementi}
\begin{concept}{Stack Operations}
Immagina l'array come una pila di piatti.
\begin{itemize}
    \item \textbf{push(x):} Aggiunge in fondo. (Veloce)
    \item \textbf{pop():} Toglie dal fondo. (Veloce)
\end{itemize}
\end{concept}
\begin{jsbox}
let numeri = [1, 2];
numeri.push(3);    // [1, 2, 3]
numeri.pop();      // [1, 2]
numeri.unshift(0); // [0, 1, 2]
\end{jsbox}

\subsection{Esercizi di Applicazione}
Per consolidare questi concetti, la lezione propone due sfide:
\begin{enumerate}
    \item \textbf{Gestione Inventario:} Creare un sistema che accumula coppie "Prodotto-Quantità" inserite dall'utente, fermandosi quando si digita "stop". Qui l'oggetto funge da database in memoria.
    \item \textbf{Media Voti:} Un programma che permette di inserire voti per vari studenti e, al termine, calcola la media per ciascuno. Questo esercizio combina probabilmente array (per i voti) e oggetti (per associare voti agli studenti).
\end{enumerate}
